\documentclass[conference]{IEEEtran}

% ── Packages ──────────────────────────────────────────────────
\usepackage[T1]{fontenc}
\usepackage[utf8]{inputenc}
\usepackage{amsmath, amssymb}
\usepackage{booktabs}
\usepackage{array}
\usepackage{multirow}
\usepackage{graphicx}
\usepackage{caption}
\usepackage{subcaption}
\usepackage[hidelinks]{hyperref}
\usepackage{cite}
\usepackage{microtype}
\usepackage{makecell}
\usepackage{url}

\graphicspath{{figures/}}

% ──────────────────────────────────────────────────────────────
\begin{document}

\title{Transfer Learning for Freshness Classification of Fruits and Vegetables:\\
A Comparative Study of Backbone Architectures and Training Conditions}

\author{
    \IEEEauthorblockN{Author Name}
    \IEEEauthorblockA{
        Department of Computer Science\\
        Institution Name\\
        City, Country\\
        email@institution.edu
    }
}

\maketitle

% ── Abstract ──────────────────────────────────────────────────
\begin{abstract}
Automated freshness classification of fruits and vegetables is a
critical task for food safety, supply chain management, and retail
quality control.
This paper presents a systematic comparative study of transfer learning
applied to this problem, evaluating three convolutional neural network
backbone architectures --- ResNet-18, EfficientNet-B0, and
EfficientNet-B2 --- under four training conditions that vary the degree
of backbone fine-tuning and the presence of a non-linear projection head.
Experiments are conducted on six publicly available datasets comprising
over 55,000 images across multiple fruit and vegetable categories,
including binary (fresh/rotten) and three-class (fresh/formalin/rotten)
annotation schemes.
We find that EfficientNet-B0 with a frozen backbone and projection head
achieves the best performance on datasets with more than 1,500 images,
reaching 100\% macro F1-score on one benchmark,
while ResNet-18 remains superior for very small datasets ($<$500 images).
A key finding is that the optimal training condition differs by backbone
family: fine-tuning the last residual block is beneficial for ResNet-18
(C4), but detrimental for EfficientNet variants, which perform best
with a completely frozen backbone (C3).
We further evaluate domain shift by testing trained models on a small
collection of 121 real-world photos taken under natural lighting
conditions, revealing a significant performance gap of up to 34.3
percentage points between studio-quality training data and
real-world evaluation images.
\end{abstract}

\begin{IEEEkeywords}
transfer learning, freshness classification, convolutional neural networks,
EfficientNet, ResNet, domain shift, food quality inspection, fine-tuning
\end{IEEEkeywords}

% ── I. Introduction ───────────────────────────────────────────
\section{Introduction}
\label{sec:intro}

Food freshness is a critical quality attribute that directly affects
consumer health, retail profitability, and supply chain efficiency.
Manual inspection of fruits and vegetables is time-consuming, subjective,
and difficult to scale to the volumes required by modern distribution
systems.
Automated computer vision systems offer a scalable and objective
alternative, capable of processing large volumes of produce at low cost.

Recent advances in deep learning, particularly convolutional neural
networks (CNNs) pre-trained on large-scale image datasets, have
demonstrated strong performance on a wide range of visual classification
tasks.
Transfer learning~\cite{pan2010transfer} leverages these pre-trained
representations to build effective classifiers from limited
domain-specific data --- an important property for food quality
applications where labeled datasets are expensive to collect.

Despite the growing body of work on food classification and freshness
detection, a systematic comparison of backbone architectures and
fine-tuning strategies under controlled conditions remains lacking.
Furthermore, the domain shift between publicly available dataset images
(typically acquired in controlled photographic conditions) and real-world
images captured by mobile devices has received limited attention.

This paper addresses both gaps through the following contributions:
\begin{itemize}
    \item A controlled comparison of three CNN backbone architectures
    (ResNet-18, EfficientNet-B0, EfficientNet-B2) and four training
    conditions (C1--C4) across six benchmark datasets totalling
    55,040 images.
    \item The identification of a consistent interaction between backbone
    family and optimal training condition: C4 (partial fine-tuning +
    projection head) for ResNet-18 vs.\ C3 (frozen backbone + projection
    head) for EfficientNet variants.
    \item A quantitative domain shift analysis showing that models
    trained on internet-sourced studio photos can lose up to 34 percentage
    points of F1-score when evaluated on real-world mobile camera images.
    \item An open-source pipeline implementing the full experimental
    workflow from data preparation to evaluation reporting.
\end{itemize}

The remainder of the paper is organized as follows.
Section~\ref{sec:related} reviews related work.
Section~\ref{sec:method} describes the methodology.
Section~\ref{sec:experiments} presents experimental results.
Section~\ref{sec:discussion} discusses key findings.
Section~\ref{sec:conclusion} concludes the paper.

% ── II. Related Work ──────────────────────────────────────────
\section{Related Work}
\label{sec:related}

\subsection{Food Quality Classification with Deep Learning}

Deep learning has been widely applied to food recognition and quality
assessment~\cite{placeholder_food}.
Several works specifically target freshness detection using CNNs trained
from scratch or with transfer learning~\cite{placeholder_fruits}.
Bangladeshi et al.\ evaluated multiple deep learning architectures for
fresh/rotten detection of fruits and vegetables, finding that pre-trained
models significantly outperform models trained from
scratch~\cite{placeholder_veg}.

\subsection{Backbone Architectures}

ResNet~\cite{he2016resnet} introduced skip connections to address the
vanishing gradient problem, enabling deeper networks that generalize well
across domains.
EfficientNet~\cite{tan2019efficientnet} proposed compound scaling ---
simultaneously adjusting network depth, width, and input resolution ---
achieving superior ImageNet accuracy with fewer parameters than
comparably-performing ResNet variants.
EfficientNet-B0 and B2 achieve 77.1\% and 80.1\% top-1 accuracy on
ImageNet respectively, compared to 69.8\% for ResNet-18, with
approximately half the parameters.

\subsection{Transfer Learning Strategies}

The most common transfer learning strategies for image classification
range from linear probing (training only the final classifier)
to full fine-tuning (updating all weights)~\cite{pan2010transfer}.
Kornblith et al.~\cite{kornblith2019better} showed that better ImageNet
features transfer better across tasks.
The intermediate strategy of fine-tuning only the later layers ---
which capture more task-specific features --- is widely adopted as a
balance between adaptation and catastrophic
forgetting~\cite{kirkpatrick2017catastrophic}.

% ── III. Methodology ──────────────────────────────────────────
\section{Methodology}
\label{sec:method}

\subsection{Datasets}

Six publicly available datasets are used in this study
(Table~\ref{tab:datasets}).
Together they cover 55,040 images across fruits, vegetables, and citrus
varieties, with binary freshness labels (fresh/rotten) in five datasets
and a three-class annotation (fresh/formalin/rotten) in one.
The FruitVision dataset is particularly relevant as the formalin class
presents a significant challenge: chemically-treated samples are
visually similar to fresh specimens, analogous to the difficulty of
detecting subtle spoilage in fermented foods.

\begin{table}[h]
\centering
\caption{Datasets used in this study. $^\dagger$Three-class annotation.}
\label{tab:datasets}
\setlength{\tabcolsep}{4pt}
\begin{tabular}{lcccc}
\toprule
\textbf{Dataset} & \textbf{ID} & \textbf{Images} & \textbf{Classes} & \textbf{Source} \\
\midrule
FruitVision$^\dagger$   & MFV & 10,154 & 3 & Mendeley \\
Fruits Classification   & MFR &  1,655 & 2 & Mendeley \\
Lemon Varieties         & MLM &  1,956 & 2 & Mendeley \\
Fruits Quality          & KFQ &    359 & 2 & Kaggle   \\
Fruits Fresh/Rotten     & KFR & 13,599 & 2 & Kaggle   \\
Fresh \& Stale          & KFS & 27,317 & 2 & Kaggle   \\
\midrule
\textbf{Total}          &     & \textbf{55,040} & & \\
\bottomrule
\end{tabular}
\end{table}

\subsection{Data Preprocessing}

All images are resized so the shortest side is 256 pixels and then
center-cropped to $224 \times 224$ for validation.
Training images undergo the following augmentation pipeline:
random resized crop to $224 \times 224$,
random horizontal flip,
$\pm 30^{\circ}$ rotation,
color jitter (brightness, contrast, saturation, hue),
and Gaussian blur (kernel size 3).
All images are normalized with ImageNet channel statistics
($\mu = [0.485, 0.456, 0.406]$, $\sigma = [0.229, 0.224, 0.225]$).
A stratified 80/20 train/validation split with fixed random seed 42
is applied to all datasets.

\subsection{Backbone Architectures}

Three pre-trained CNN backbones are evaluated, all initialized with
ImageNet weights from the \texttt{torchvision} library:

\begin{itemize}
    \item \textbf{ResNet-18}~\cite{he2016resnet}: 11.2M parameters,
    512-dimensional feature output from global average pooling.
    \item \textbf{EfficientNet-B0}~\cite{tan2019efficientnet}: 5.3M
    parameters, 1280-dimensional feature output.
    \item \textbf{EfficientNet-B2}~\cite{tan2019efficientnet}: 9.1M
    parameters, 1408-dimensional feature output.
\end{itemize}

\subsection{Training Conditions}

Four training conditions form an ablation study over backbone
fine-tuning and classification head architecture (Table~\ref{tab:conditions}).

\begin{table}[h]
\centering
\caption{Training conditions.}
\label{tab:conditions}
\begin{tabular}{lp{1.8cm}p{1.8cm}p{2.2cm}}
\toprule
\textbf{ID} & \textbf{Backbone} & \textbf{Head} & \textbf{Description} \\
\midrule
C1 & Frozen      & Linear only         & Linear probe \\
C2 & Layer4 free & Linear only         & Partial fine-tuning \\
C3 & Frozen      & Proj.\ + Linear     & Full head, frozen backbone \\
C4 & Layer4 free & Proj.\ + Linear     & Full head + fine-tuning \\
\bottomrule
\end{tabular}
\end{table}

In conditions C3 and C4, a two-layer projection head maps the backbone
output to a 128-dimensional embedding before classification:
$d_\text{out} \xrightarrow{\text{FC}} 256 \xrightarrow{\text{ReLU}}
\xrightarrow{\text{FC}} 128 \xrightarrow{\text{ReLU}} C$,
where $C$ is the number of classes.
In condition C4, \texttt{layer4} of the backbone is unfrozen and trained
with a reduced learning rate $\text{lr}_\text{backbone} = 10^{-5}$,
compared to $\text{lr} = 10^{-3}$ for the head, to mitigate catastrophic
forgetting~\cite{kirkpatrick2017catastrophic}.

All models are trained for 60 epochs using the Adam optimizer with
cosine annealing learning rate schedule and weight decay $10^{-4}$.
Batch size is 32.

\subsection{Evaluation Metrics}

Given the potential class imbalance in food datasets, accuracy alone is
insufficient. Let $TP_c$, $FP_c$, $FN_c$ denote true positives, false
positives, and false negatives for class $c$, and $C$ the number of
classes.
We report:
\begin{align}
    \overline{F1} &= \frac{1}{C}\sum_{c=1}^{C}
    \frac{2\,TP_c}{2\,TP_c + FP_c + FN_c} \\
    \text{MCC} &= \frac{TP \cdot TN - FP \cdot FN}
    {\sqrt{(TP+FP)(TP+FN)(TN+FP)(TN+FN)}}
\end{align}
Additionally, we report Area Under the ROC Curve (AUC-ROC) and
per-class F1, precision, and recall.
For food safety applications, recall on the \textit{rotten} class is
particularly critical: a false negative (rotten food classified as fresh)
poses a direct health risk.

% ── IV. Experiments ───────────────────────────────────────────
\section{Experiments}
\label{sec:experiments}

\subsection{Phase 1: ResNet-18 Baseline}

Table~\ref{tab:r18} reports the best macro F1-score achieved by
ResNet-18 for each dataset and condition (state mode, 60 epochs).
Condition C4 achieves the best result in four of six datasets;
C2 wins on the two datasets where the number of training samples
is smallest (KFQ: 287 train samples; MFV: 8,123 train samples
but three-class complexity).

\begin{table}[h]
\centering
\caption{Best macro F1-score (\%) --- ResNet-18, state mode.
Bold: best condition per dataset.}
\label{tab:r18}
\setlength{\tabcolsep}{5pt}
\begin{tabular}{lccccl}
\toprule
\textbf{Dataset} & \textbf{C1} & \textbf{C2} & \textbf{C3} & \textbf{C4} & \textbf{Best} \\
\midrule
KFR & 97.6 & 99.8 & 99.4 & \textbf{99.8} & C2/C4 \\
KFS & 95.0 & 98.9 & 98.0 & \textbf{99.0} & C4 \\
MLM & 96.2 & 98.0 & 96.2 & \textbf{98.5} & C4 \\
MFR & 94.3 & 95.5 & 95.2 & \textbf{96.4} & C4 \\
MFV & 78.3 & \textbf{91.0} & 86.7 & 90.6 & C2 \\
KFQ & 86.1 & \textbf{91.6} & 86.0 & 88.9 & C2 \\
\bottomrule
\end{tabular}
\end{table}

\subsection{Phase 2: Backbone Comparison}

Table~\ref{tab:backbone} compares the best macro F1-score per dataset
across the three backbone architectures (state mode, best condition).
EfficientNet-B0 with condition C3 achieves the best result in three
datasets; EfficientNet-B2 wins on MFV (the hardest, three-class dataset)
and matches EB0 on KFS; ResNet-18 wins only on KFQ (359 images),
the smallest dataset in the benchmark.
Notably, EB2 does not consistently outperform EB0 despite having
more parameters, suggesting diminishing returns from increased
model capacity for this task.

\begin{table}[h]
\centering
\caption{Best macro F1-score (\%) by backbone architecture, state mode.
Bold: overall best per dataset.}
\label{tab:backbone}
\setlength{\tabcolsep}{4pt}
\begin{tabular}{lccc ccc}
\toprule
\multirow{2}{*}{\textbf{Dataset}}
    & \multicolumn{2}{c}{\textbf{ResNet-18}}
    & \multicolumn{2}{c}{\textbf{EffNet-B0}}
    & \multicolumn{2}{c}{\textbf{EffNet-B2}} \\
\cmidrule(lr){2-3}\cmidrule(lr){4-5}\cmidrule(lr){6-7}
    & F1 & Cond & F1 & Cond & F1 & Cond \\
\midrule
KFR & 99.8 & C4 & \textbf{100.0} & C3 & 99.9 & C3 \\
KFS & 99.0 & C4 & \textbf{99.4} & C3 & 99.0 & C3 \\
MLM & 98.5 & C4 & \textbf{98.7} & C3 & 98.5 & C3 \\
MFR & 96.4 & C4 & \textbf{97.0} & C3 & 95.8 & C4 \\
MFV & 91.0 & C2 & 91.6 & C3 & \textbf{92.3} & C3 \\
KFQ & \textbf{91.6} & C2 & 86.1 & C4 & 88.9 & C3 \\
\bottomrule
\end{tabular}
\end{table}

\subsection{Domain Shift Analysis}

To measure generalization to real-world conditions, trained models are
evaluated on \textit{own\_dataset}: 121 strawberry photos taken with a
mobile phone camera under natural indoor lighting.
None of the training datasets contain photos of strawberries taken in
these conditions.
Table~\ref{tab:domain} reports the top results sorted by evaluation
F1-score, with the domain shift (Drop = F1$_\text{train}$ $-$
F1$_\text{eval}$) for each model.

\begin{table}[h]
\centering
\caption{Cross-dataset evaluation on \textit{own\_dataset} (121 real
strawberry photos, state mode). Top results and notable extremes.
Drop = F1$_\text{train}$ $-$ F1$_\text{eval}$.}
\label{tab:domain}
\setlength{\tabcolsep}{4pt}
\begin{tabular}{lcccc}
\toprule
\textbf{Model ID} & \textbf{F1$_\text{train}$} & \textbf{F1$_\text{eval}$}
    & \textbf{Recall$_\text{eval}$} & \textbf{Drop} \\
\midrule
MLM-R18-C4  & 98.5 & \textbf{91.7} & 92.3 & 6.8  \\
MLM-EB0-C4  & 95.4 & \textbf{91.7} & 92.3 & 3.7  \\
MFR-EB2-C3  & ---  & \textbf{91.6} & 91.6 & ---  \\
MLM-R18-C2  & 98.0 & 87.5          & 88.5 & 10.5 \\
MFR-EB0-C3  & 97.0 & 87.3          & 87.1 & 9.7  \\
KFS-R18-C4  & 99.0 & 87.3          & 87.1 & 11.7 \\
KFR-R18-C4  & 99.8 & 70.8          & 71.0 & 29.0 \\
KFR-EB0-C3  & 100.0 & 65.7         & 65.7 & 34.3 \\
MLM-EB2-C4  & ---  & 43.4          & 47.9 & ---  \\
\bottomrule
\end{tabular}
\end{table}

% ── V. Results and Discussion ─────────────────────────────────
\section{Results and Discussion}
\label{sec:discussion}

\textbf{Finding 1: C3 is optimal for EfficientNet; C4 for ResNet-18.}
In four of five completed datasets, EfficientNet-B0 achieves its best
performance under C3 (frozen backbone + projection head), while ResNet-18
consistently benefits from partial fine-tuning (C4).
We attribute this to the richer pre-trained features of EfficientNet:
its 1280-dimensional output already captures discriminative cues for
freshness without backbone adaptation, whereas ResNet-18's 512-dimensional
output benefits from specialized high-level feature adaptation.
Fine-tuning EfficientNet's backbone with the available data sizes
appears to introduce overfitting rather than beneficial specialization.

\textbf{Finding 2: Dataset size determines the winning backbone.}
Below approximately 500 training samples, ResNet-18 outperforms
EfficientNet-B0 by up to 5.5 percentage points (KFQ dataset: 91.6\%
vs.\ 86.1\%).
Above 1,500 training samples, EfficientNet-B0 consistently matches
or surpasses ResNet-18.
This is consistent with the larger head parameters of EfficientNet
requiring more data to generalize: a $1280 \to 256 \to 128 \to C$
projection head has approximately 3$\times$ more parameters than the
equivalent ResNet-18 head.

\textbf{Finding 3: EfficientNet-B0 achieves perfect classification on KFR.}
On the Fruits Fresh/Rotten dataset (apple, banana, orange),
EfficientNet-B0 under C3 achieves 100\% macro F1-score, MCC = 1.0,
and AUC-ROC = 1.0.
This represents a perfect separation between fresh and rotten classes,
suggesting that the visual differences in this dataset are highly
discriminative and well-captured by EfficientNet's pre-trained features
without backbone adaptation.

\textbf{Finding 4: Formalin represents a fundamental classification challenge.}
FruitVision is the only dataset where no backbone exceeds 92.3\% F1-score,
despite having 10,154 images.
The formalin class is chemically treated to appear visually similar to
fresh samples, creating a subtle spectral shift that is difficult to
detect from RGB images alone.
EfficientNet-B2 achieves the best result (92.3\%), suggesting that
larger, more expressive backbones provide marginal improvements on
hard discrimination tasks.
This finding is directly relevant to the detection of subtle spoilage
in fermented foods.

\textbf{Finding 5: High training F1 does not imply good generalization.}
The Fruits Fresh/Rotten model that achieves 100\% F1 on its own training
domain (KFR-EB0-C3) performs worst in cross-dataset evaluation,
dropping to 65.7\% on real photos (drop of 34.3 points).
Conversely, the Lemon Varieties model (MLM-R18-C4, 98.5\% F1) achieves
the highest evaluation performance (91.7\%, drop of only 6.8 points).
We attribute this inversion to the diverse photographic conditions in
the Lemon Varieties dataset, which expose the model to a wider range
of lighting and background variation during training.

\textbf{Finding 6: Domain shift is quantifiable and large.}
Across 60 evaluated models, the mean domain shift is 16.0 percentage
points, with a range from 3.7 to 34.3 points.
Models trained on datasets with more uniform photographic conditions
(KFR: studio photos with white backgrounds) exhibit the largest shift.
This strongly motivates the collection of domain-specific datasets
that match real deployment conditions when high-accuracy production
systems are required.

\textbf{Finding 7: EfficientNet-B2 exhibits high variance in generalization.}
While EfficientNet-B2 achieves competitive or superior training performance,
its cross-dataset behavior is highly inconsistent.
MFR-EB2-C3 achieves 91.6\% F1 on \textit{own\_dataset} (top result),
while MLM-EB2-C4 drops to 43.4\% (worst result across all models).
This 48-point range --- compared to a 22-point range for ResNet-18
and a 26-point range for EfficientNet-B0 --- suggests that EB2's
larger feature dimensionality (1408) amplifies both the benefits
of good training data alignment and the costs of misalignment.
The condition C4 appears particularly problematic for EB2:
all EB2 C4 models in the bottom half of Table~\ref{tab:domain}
belong to this condition.

\textbf{Finding 8: MFR-EB2-C1 (frozen backbone, linear only) ties with top models.}
MFR-EB2-C3 (frozen + projection head) and MFR-EB2-C1 (frozen + linear only)
both achieve 91.6\% F1 on \textit{own\_dataset}.
This suggests that for EB2, even the simplest linear probe already
captures sufficient discriminative information from the pre-trained
features for this specific cross-domain evaluation, and that the
projection head provides no additional benefit in this setting.
The result further reinforces the finding that fine-tuning is not
necessary --- and may be harmful --- for larger EfficientNet models.

% ── VI. Conclusion ────────────────────────────────────────────
\section{Conclusion}
\label{sec:conclusion}

This paper presented a systematic comparison of transfer learning
strategies for fruit and vegetable freshness classification.
Three backbone architectures (ResNet-18, EfficientNet-B0, EfficientNet-B2)
and four training conditions were evaluated across six benchmark datasets
totalling 55,040 images.
Cross-dataset generalization was measured against 121 real-world photos
across 60 model evaluations.

The key findings are:
(1) EfficientNet-B0 with a frozen backbone and projection head (C3)
achieves the best training performance on medium-to-large datasets,
including a perfect 100\% F1-score on one benchmark;
(2) ResNet-18 with layer4 fine-tuning (C4) remains the best choice for
datasets with fewer than 500 training samples;
(3) the optimal condition differs by backbone family --- C4 for ResNet-18,
C3 for EfficientNet variants --- suggesting that richer pre-trained
features reduce the need for backbone fine-tuning;
(4) domain shift ranges from 3.7 to 34.3 percentage points across
60 evaluations, with a mean of 16.0 points;
(5) EfficientNet-B2 shows high variance in generalization (range:
43.4\%--91.6\%), making it less reliable than EB0 for
cross-domain deployment despite superior training performance on
challenging datasets.

Future work will extend this pipeline to Korean food categories
(kimchi, rice cake, fermented side dishes), where the visual similarity
between fresh and fermented states is directly analogous to the
formalin challenge identified in this study.

% ── References ────────────────────────────────────────────────
\bibliographystyle{IEEEtran}
\bibliography{references}

\end{document}
